\section{Bonus}

\subsection{Dokumentasi API}
% TODO: ExDoc, tautan hasil hosting.

\subsection{Signed SBOM}

SBOM dibuat dari \texttt{mix.lock} dan berformat JSON dengan schema \texttt{minisign-sbom/v1}.
Isinya \textit{subject}, yaitu nama, ukuran, dan SHA-256 artefak rilis yang dirujuk, serta
daftar \textit{components} yang diurutkan menurut nama. Setiap komponen memuat nama, versi,
\textit{purl} (\texttt{pkg:hex/<nama>@<versi>}), dan checksum tarball SHA-256 dari
\texttt{mix.lock}. Pembentukan bersifat deterministik: masukan yang sama selalu menghasilkan
byte yang sama.

SBOM ditandatangani sebagai artefak terpisah memakai alur penandatanganan yang sama dengan
rilis, sehingga bundle SBOM memakai format bundle yang sama (Tabel 2.4). Verifikasi
memeriksa dua bundle dan mencocokkan subject dengan artefak rilis.

\begin{table}[H]
\centering
\begin{tabular}{@{}llp{7.5cm}@{}}
\toprule
Subcommand & Parameter & Keluaran \\
\midrule
\texttt{sbom generate} & \texttt{-file}, \texttt{-lock}, \texttt{-out} &
SBOM JSON pada \texttt{-out}. Tidak menimpa berkas yang sudah ada. \\
\texttt{sbom sign} & \texttt{-sbom}, \texttt{-package}, \texttt{-version}, \texttt{-key}, \texttt{-out} &
Bundle SBOM pada \texttt{-out}. \\
\texttt{sbom verify} & \texttt{-file}, \texttt{-bundle}, \texttt{-sbom}, \texttt{-sbom-bundle}, \texttt{-trust-store},
\texttt{-package}, \texttt{-version}, \texttt{-json} (opsional) &
Hasil verifikasi sesuai Tabel 2.9. Exit code sesuai Tabel 2.10. \\
\bottomrule
\end{tabular}
\caption{Subcommand tambahan untuk SBOM.}
\end{table}

Verifikasi SBOM berjalan berurutan dan berhenti pada kegagalan pertama. Langkah pertama
memverifikasi bundle SBOM dengan alur tujuh langkah yang sama. Langkah kedua memverifikasi
bundle rilis. Langkah ketiga mencocokkan nama, ukuran, dan SHA-256 pada subject dengan
artefak rilis. Kegagalan langkah ketiga dilaporkan sebagai ARTIFACT\_MISMATCH, dan SBOM yang
formatnya tidak sesuai dilaporkan sebagai FORMAT\_ERROR.

\paragraph{Pengujian.}
Pengujian SBOM dijalankan lewat CLI pada berkas \texttt{test/sbom\_test.exs} (14 tes, seluruhnya
lolos) dan diuji ulang di dalam image dengan UID acak dan \texttt{--network none}.

\begin{table}[H]
\centering
\begin{tabular}{@{}lll@{}}
\toprule
Kasus uji & Hasil yang diharapkan & Hasil \\
\midrule
SBOM valid dan saling merujuk & OK, exit 0 & OK, exit 0 \\
Isi SBOM diubah & ARTIFACT\_MISMATCH, exit 1 & ARTIFACT\_MISMATCH, exit 1 \\
Artefak rilis diganti isinya & ARTIFACT\_MISMATCH, exit 1 & ARTIFACT\_MISMATCH, exit 1 \\
SBOM milik artefak lain, bundle sah & ARTIFACT\_MISMATCH, exit 1 & ARTIFACT\_MISMATCH, exit 1 \\
Bundle SBOM dari kunci tak tepercaya & KEY\_NOT\_TRUSTED, exit 1 & KEY\_NOT\_TRUSTED, exit 1 \\
SBOM rusak namun ditandatangani & FORMAT\_ERROR, exit 1 & FORMAT\_ERROR, exit 1 \\
Berkas SBOM tidak ada & INPUT\_ERROR, exit 2 & INPUT\_ERROR, exit 2 \\
\bottomrule
\end{tabular}
\caption{Hasil pengujian SBOM.}
\end{table}

\subsection{Video}
% TODO: \videourl
